\documentclass[12pt,a4paper]{article}
\usepackage[utf8]{inputenc}
\usepackage{amsmath,amssymb,amsthm}
\usepackage{geometry}
\usepackage{graphicx}
\usepackage{hyperref}
\usepackage{enumitem}
\usepackage{xcolor}
\usepackage{tcolorbox}
\tcbuselibrary{breakable,skins}
\usepackage{mathtools}
\usepackage{fancyhdr}
\usepackage{tikz}
\usetikzlibrary{arrows.meta,positioning,calc}

\geometry{margin=1in}

\pagestyle{fancy}
\fancyhf{}
\rhead{Matrix Operations in ML \& DL}
\lhead{Detailed Notes}
\cfoot{\thepage}

\title{\textbf{Linear Transformations and Vector Operations \\ for Machine Learning and Deep Learning}}
\author{Akshay Redekar}
\date{\today}

\theoremstyle{definition}
\newtheorem{definition}{Definition}[section]
\newtheorem{theorem}{Theorem}[section]
\newtheorem{example}{Example}[section]
\newtheorem{property}{Property}[section]

\definecolor{highlight}{RGB}{255,255,200}

% Add page border using tikz
\AddToHook{shipout/background}{%
    \begin{tikzpicture}[remember picture, overlay]
        \draw[line width=0.5pt]
            ($(current page.north west) + (0.2in,-0.2in)$)
            rectangle
            ($(current page.south east) + (-0.2in,0.2in)$);
    \end{tikzpicture}
}

\begin{document}

\maketitle
\tableofcontents
\newpage

\section{Linear Transformation of Non-square Matrix}

\subsection{Introduction and Motivation}

In machine learning and deep learning, we frequently encounter non-square matrices because the number of features doesn't usually match the number of data points. Understanding how these matrices transform vectors is fundamental to many ML algorithms.

\begin{definition}[Non-square Matrix]
A non-square matrix $\mathbf{A} \in \mathbb{R}^{m \times n}$ where $m \neq n$ represents a linear transformation that maps vectors from $\mathbb{R}^n$ to $\mathbb{R}^m$:
\begin{equation}
\mathbf{A}: \mathbb{R}^n \to \mathbb{R}^m
\end{equation}

\begin{itemize}
    \item If $m > n$: The transformation maps from lower to higher dimension (expansion)
    \item If $m < n$: The transformation maps from higher to lower dimension (reduction)
\end{itemize}
\end{definition}

\begin{tcolorbox}[enhanced,breakable,colback=blue!5!white,colframe=blue!75!black,title=Why Non-square Matrices in ML/DL?]
\textbf{Common Scenarios:}
\begin{enumerate}
    \item \textbf{Dataset Matrix:} $\mathbf{X} \in \mathbb{R}^{n \times d}$ where $n$ samples $\neq$ $d$ features
    \item \textbf{Neural Network Layers:} Input dimension often differs from output dimension
    \item \textbf{Dimensionality Reduction:} PCA, autoencoders map from high to low dimensions
    \item \textbf{Feature Extraction:} Convolutional layers change spatial dimensions
    \item \textbf{Embedding Layers:} Map discrete tokens to continuous vectors (different dimensions)
\end{enumerate}
\end{tcolorbox}

\subsection{Mathematical Properties}

\begin{theorem}[Properties of Non-square Linear Transformations]
For matrix $\mathbf{A} \in \mathbb{R}^{m \times n}$ with $m \neq n$:

\begin{enumerate}
    \item \textbf{Linearity:} $\mathbf{A}(c_1\mathbf{x}_1 + c_2\mathbf{x}_2) = c_1\mathbf{A}\mathbf{x}_1 + c_2\mathbf{A}\mathbf{x}_2$
    
    \item \textbf{Rank:} $\text{rank}(\mathbf{A}) \leq \min(m, n)$
    
    \item \textbf{Not Invertible:} Cannot have a true inverse since $m \neq n$, but may have:
    \begin{itemize}
        \item Left inverse if $m > n$ and full column rank
        \item Right inverse if $m < n$ and full row rank
        \item Pseudo-inverse in general case
    \end{itemize}
    
    \item \textbf{Null Space:} If $m < n$, the null space is non-trivial: $\dim(\mathcal{N}(\mathbf{A})) = n - \text{rank}(\mathbf{A}) > 0$
\end{enumerate}
\end{theorem}

\subsection{Types of Non-square Transformations}

\subsubsection{Expansion: From Lower to Higher Dimension}

\begin{example}[Dimension Expansion $\mathbb{R}^2 \to \mathbb{R}^3$]
Consider transformation matrix:
\begin{equation}
\mathbf{A} = \begin{bmatrix}
1 & 0 \\
0 & 1 \\
1 & 1
\end{bmatrix} \in \mathbb{R}^{3 \times 2}
\end{equation}

\textbf{Step 1: Apply to input vector}
\begin{equation}
\mathbf{x} = \begin{bmatrix} 2 \\ 3 \end{bmatrix} \in \mathbb{R}^2
\end{equation}

\textbf{Step 2: Compute transformation}
\begin{align}
\mathbf{y} = \mathbf{A}\mathbf{x} &= \begin{bmatrix}
1 & 0 \\
0 & 1 \\
1 & 1
\end{bmatrix} \begin{bmatrix} 2 \\ 3 \end{bmatrix} \\
&= \begin{bmatrix}
1(2) + 0(3) \\
0(2) + 1(3) \\
1(2) + 1(3)
\end{bmatrix} = \begin{bmatrix} 2 \\ 3 \\ 5 \end{bmatrix} \in \mathbb{R}^3
\end{align}

\textbf{Interpretation:}
\begin{itemize}
    \item First component: keeps $x_1$
    \item Second component: keeps $x_2$
    \item Third component: sum $x_1 + x_2$
\end{itemize}

\textbf{Geometric view:} The 2D vector is embedded into 3D space along a plane.

\textbf{Range:} All outputs lie in a 2D plane within $\mathbb{R}^3$:
\begin{equation}
\mathcal{C}(\mathbf{A}) = \text{span}\left\{\begin{bmatrix} 1 \\ 0 \\ 1 \end{bmatrix}, \begin{bmatrix} 0 \\ 1 \\ 1 \end{bmatrix}\right\}
\end{equation}
\end{example}

\subsubsection{Reduction: From Higher to Lower Dimension}

\begin{example}[Dimension Reduction $\mathbb{R}^3 \to \mathbb{R}^2$]
Consider transformation matrix:
\begin{equation}
\mathbf{A} = \begin{bmatrix}
1 & 0 & 1 \\
0 & 1 & 1
\end{bmatrix} \in \mathbb{R}^{2 \times 3}
\end{equation}

\textbf{Step 1: Apply to input vector}
\begin{equation}
\mathbf{x} = \begin{bmatrix} 2 \\ 3 \\ 1 \end{bmatrix} \in \mathbb{R}^3
\end{equation}

\textbf{Step 2: Compute transformation}
\begin{align}
\mathbf{y} = \mathbf{A}\mathbf{x} &= \begin{bmatrix}
1 & 0 & 1 \\
0 & 1 & 1
\end{bmatrix} \begin{bmatrix} 2 \\ 3 \\ 1 \end{bmatrix} \\
&= \begin{bmatrix}
1(2) + 0(3) + 1(1) \\
0(2) + 1(3) + 1(1)
\end{bmatrix} = \begin{bmatrix} 3 \\ 4 \end{bmatrix} \in \mathbb{R}^2
\end{align}

\textbf{Interpretation:}
\begin{itemize}
    \item First component: $x_1 + x_3$
    \item Second component: $x_2 + x_3$
    \item Information from 3D is projected to 2D
\end{itemize}

\textbf{Null Space:} Find vectors mapped to zero.

Solve $\mathbf{A}\mathbf{x} = \mathbf{0}$:
\begin{equation}
\begin{bmatrix}
1 & 0 & 1 \\
0 & 1 & 1
\end{bmatrix} \begin{bmatrix} x_1 \\ x_2 \\ x_3 \end{bmatrix} = \begin{bmatrix} 0 \\ 0 \end{bmatrix}
\end{equation}

From equations:
\begin{align}
x_1 + x_3 &= 0 \implies x_1 = -x_3 \\
x_2 + x_3 &= 0 \implies x_2 = -x_3
\end{align}

Let $x_3 = t$:
\begin{equation}
\mathbf{x} = t\begin{bmatrix} -1 \\ -1 \\ 1 \end{bmatrix}
\end{equation}

\textbf{Null space basis:}
\begin{equation}
\mathcal{N}(\mathbf{A}) = \text{span}\left\{\begin{bmatrix} -1 \\ -1 \\ 1 \end{bmatrix}\right\}
\end{equation}

\textbf{Dimension analysis:}
\begin{itemize}
    \item $\text{rank}(\mathbf{A}) = 2$
    \item $\dim(\mathcal{N}(\mathbf{A})) = 3 - 2 = 1$ ✓
\end{itemize}

This shows information loss: any vector along the null space direction is mapped to zero.
\end{example}

\subsection{Left and Right Inverses}

\begin{definition}[Left Inverse]
For $\mathbf{A} \in \mathbb{R}^{m \times n}$ with $m > n$ and full column rank, the left inverse is:
\begin{equation}
\mathbf{A}_L^{-1} = (\mathbf{A}^T\mathbf{A})^{-1}\mathbf{A}^T
\end{equation}

Property: $\mathbf{A}_L^{-1}\mathbf{A} = \mathbf{I}_n$ (but $\mathbf{A}\mathbf{A}_L^{-1} \neq \mathbf{I}_m$)
\end{definition}

\begin{definition}[Right Inverse]
For $\mathbf{A} \in \mathbb{R}^{m \times n}$ with $m < n$ and full row rank, the right inverse is:
\begin{equation}
\mathbf{A}_R^{-1} = \mathbf{A}^T(\mathbf{A}\mathbf{A}^T)^{-1}
\end{equation}

Property: $\mathbf{A}\mathbf{A}_R^{-1} = \mathbf{I}_m$ (but $\mathbf{A}_R^{-1}\mathbf{A} \neq \mathbf{I}_n$)
\end{definition}

\begin{example}[Computing Left Inverse]
For:
\begin{equation}
\mathbf{A} = \begin{bmatrix}
1 & 0 \\
0 & 1 \\
1 & 1
\end{bmatrix} \in \mathbb{R}^{3 \times 2}
\end{equation}

\textbf{Step 1: Compute $\mathbf{A}^T\mathbf{A}$}
\begin{equation}
\mathbf{A}^T\mathbf{A} = \begin{bmatrix}
1 & 0 & 1 \\
0 & 1 & 1
\end{bmatrix} \begin{bmatrix}
1 & 0 \\
0 & 1 \\
1 & 1
\end{bmatrix} = \begin{bmatrix}
2 & 1 \\
1 & 2
\end{bmatrix}
\end{equation}

\textbf{Step 2: Invert $\mathbf{A}^T\mathbf{A}$}
\begin{equation}
\det(\mathbf{A}^T\mathbf{A}) = 4 - 1 = 3
\end{equation}
\begin{equation}
(\mathbf{A}^T\mathbf{A})^{-1} = \frac{1}{3}\begin{bmatrix}
2 & -1 \\
-1 & 2
\end{bmatrix}
\end{equation}

\textbf{Step 3: Compute left inverse}
\begin{align}
\mathbf{A}_L^{-1} &= (\mathbf{A}^T\mathbf{A})^{-1}\mathbf{A}^T \\
&= \frac{1}{3}\begin{bmatrix}
2 & -1 \\
-1 & 2
\end{bmatrix} \begin{bmatrix}
1 & 0 & 1 \\
0 & 1 & 1
\end{bmatrix} \\
&= \frac{1}{3}\begin{bmatrix}
2 & -1 & 1 \\
-1 & 2 & 1
\end{bmatrix}
\end{equation}

\textbf{Verification:}
\begin{align}
\mathbf{A}_L^{-1}\mathbf{A} &= \frac{1}{3}\begin{bmatrix}
2 & -1 & 1 \\
-1 & 2 & 1
\end{bmatrix} \begin{bmatrix}
1 & 0 \\
0 & 1 \\
1 & 1
\end{bmatrix} \\
&= \frac{1}{3}\begin{bmatrix}
3 & 0 \\
0 & 3
\end{bmatrix} = \begin{bmatrix}
1 & 0 \\
0 & 1
\end{bmatrix} = \mathbf{I}_2 \checkmark
\end{align}
\end{example}

\subsection{Applications in Machine Learning}

\begin{tcolorbox}[enhanced,breakable,colback=green!5!white,colframe=green!75!black,title=ML/DL Applications of Non-square Transformations]

\textbf{1. Dimensionality Reduction (PCA):}

Transform high-dimensional data to lower dimensions:
\begin{equation}
\mathbf{Z} = \mathbf{X}\mathbf{W} \quad \text{where } \mathbf{W} \in \mathbb{R}^{d \times k}, k < d
\end{equation}

Example: Project 1000-dimensional text features to 50 dimensions.

\textbf{2. Neural Network Layers:}

Each layer performs non-square transformation:
\begin{equation}
\mathbf{h}^{(l)} = \sigma(\mathbf{W}^{(l)}\mathbf{h}^{(l-1)} + \mathbf{b}^{(l)})
\end{equation}

where $\mathbf{W}^{(l)} \in \mathbb{R}^{n_l \times n_{l-1}}$ (often $n_l \neq n_{l-1}$)

\textbf{3. Feature Engineering:}

Create new features through linear combinations:
\begin{equation}
\mathbf{X}_{new} = \mathbf{X}_{old}\mathbf{T} \quad \text{where } \mathbf{T} \in \mathbb{R}^{d_{old} \times d_{new}}
\end{equation}

\textbf{4. Autoencoders:}

\begin{itemize}
    \item Encoder: $\mathbf{z} = f(\mathbf{W}_e\mathbf{x})$ where $\mathbf{W}_e \in \mathbb{R}^{k \times d}$, $k < d$
    \item Decoder: $\hat{\mathbf{x}} = g(\mathbf{W}_d\mathbf{z})$ where $\mathbf{W}_d \in \mathbb{R}^{d \times k}$
\end{itemize}

\textbf{5. Word Embeddings (NLP):}

Map vocabulary to dense vectors:
\begin{equation}
\mathbf{E} \in \mathbb{R}^{V \times d} \quad \text{where } V \text{ is vocab size, } d \text{ is embedding dim}
\end{equation}

Typically $V \gg d$ (e.g., 50000 → 300)

\textbf{6. Convolutional Layers:}

Flatten spatial dimensions while increasing channels:
\begin{equation}
\text{Input: } 32 \times 32 \times 3 \to \text{Output: } 16 \times 16 \times 64
\end{equation}
\end{tcolorbox}

\begin{example}[Complete ML Example: PCA Dimensionality Reduction]
Dataset with 4 samples, 3 features:
\begin{equation}
\mathbf{X} = \begin{bmatrix}
1 & 2 & 3 \\
2 & 3 & 5 \\
3 & 5 & 7 \\
4 & 6 & 9
\end{bmatrix} \in \mathbb{R}^{4 \times 3}
\end{equation}

\textbf{Goal:} Reduce from 3 to 2 dimensions.

\textbf{Step 1: Center the data}
\begin{equation}
\boldsymbol{\mu} = \begin{bmatrix} 2.5 \\ 4 \\ 6 \end{bmatrix}
\end{equation}
\begin{equation}
\mathbf{X}_{centered} = \begin{bmatrix}
-1.5 & -2 & -3 \\
-0.5 & -1 & -1 \\
0.5 & 1 & 1 \\
1.5 & 2 & 3
\end{bmatrix}
\end{equation}

\textbf{Step 2: Compute covariance matrix}
\begin{equation}
\mathbf{C} = \frac{1}{3}\mathbf{X}_{centered}^T\mathbf{X}_{centered} = \begin{bmatrix}
1.667 & 2.167 & 3.167 \\
2.167 & 2.917 & 4.250 \\
3.167 & 4.250 & 6.667
\end{bmatrix}
\end{equation}

\textbf{Step 3: Eigendecomposition}

Eigenvalues: $\lambda_1 = 11.01$, $\lambda_2 = 0.24$, $\lambda_3 \approx 0$

Top 2 eigenvectors (principal components):
\begin{equation}
\mathbf{W} = \begin{bmatrix}
0.425 & -0.808 \\
0.566 & -0.118 \\
0.706 & 0.578
\end{bmatrix} \in \mathbb{R}^{3 \times 2}
\end{equation}

\textbf{Step 4: Project data (non-square transformation)}
\begin{align}
\mathbf{Z} &= \mathbf{X}_{centered}\mathbf{W} \\
&= \begin{bmatrix}
-1.5 & -2 & -3 \\
-0.5 & -1 & -1 \\
0.5 & 1 & 1 \\
1.5 & 2 & 3
\end{bmatrix} \begin{bmatrix}
0.425 & -0.808 \\
0.566 & -0.118 \\
0.706 & 0.578
\end{bmatrix} \\
&\approx \begin{bmatrix}
-4.586 & 0.321 \\
-1.531 & 0.069 \\
1.531 & -0.069 \\
4.586 & -0.321
\end{bmatrix} \in \mathbb{R}^{4 \times 2}
\end{align}

\textbf{Result:} Data reduced from 3D to 2D!

\textbf{Variance explained:}
\begin{equation}
\frac{\lambda_1 + \lambda_2}{\lambda_1 + \lambda_2 + \lambda_3} = \frac{11.25}{11.25} \approx 99.8\%
\end{equation}

\textbf{Reconstruction (approximate):}
\begin{equation}
\mathbf{X}_{reconstructed} = \mathbf{Z}\mathbf{W}^T + \boldsymbol{\mu}
\end{equation}

This demonstrates how non-square transformation $\mathbf{W} \in \mathbb{R}^{3 \times 2}$ reduces dimensionality while preserving 99.8\% of variance.
\end{example}

\begin{example}[Deep Learning Example: Neural Network Layer]
Consider a fully connected layer:

\textbf{Input:} $\mathbf{x} \in \mathbb{R}^{784}$ (flattened 28×28 image)
\textbf{Output:} $\mathbf{h} \in \mathbb{R}^{128}$ (hidden layer)

\textbf{Transformation:}
\begin{equation}
\mathbf{h} = \text{ReLU}(\mathbf{W}\mathbf{x} + \mathbf{b})
\end{equation}

where $\mathbf{W} \in \mathbb{R}^{128 \times 784}$ (non-square!)

\textbf{Forward pass example (simplified):}

Let's use smaller dimensions for illustration:
\begin{equation}
\mathbf{W} = \begin{bmatrix}
0.1 & 0.2 & 0.3 & 0.1 \\
0.2 & 0.1 & 0.2 & 0.3 \\
0.3 & 0.3 & 0.1 & 0.2
\end{bmatrix} \in \mathbb{R}^{3 \times 4}, \quad
\mathbf{b} = \begin{bmatrix} 0.1 \\ 0.2 \\ 0.1 \end{bmatrix}
\end{equation}

\textbf{Input:}
\begin{equation}
\mathbf{x} = \begin{bmatrix} 1.0 \\ 2.0 \\ 1.5 \\ 0.5 \end{bmatrix} \in \mathbb{R}^4
\end{equation}

\textbf{Step 1: Linear transformation}
\begin{align}
\mathbf{z} = \mathbf{W}\mathbf{x} + \mathbf{b} &= \begin{bmatrix}
0.1 & 0.2 & 0.3 & 0.1 \\
0.2 & 0.1 & 0.2 & 0.3 \\
0.3 & 0.3 & 0.1 & 0.2
\end{bmatrix} \begin{bmatrix} 1.0 \\ 2.0 \\ 1.5 \\ 0.5 \end{bmatrix} + \begin{bmatrix} 0.1 \\ 0.2 \\ 0.1 \end{bmatrix} \\
&= \begin{bmatrix}
0.1 + 0.4 + 0.45 + 0.05 \\
0.2 + 0.2 + 0.3 + 0.15 \\
0.3 + 0.6 + 0.15 + 0.1
\end{bmatrix} + \begin{bmatrix} 0.1 \\ 0.2 \\ 0.1 \end{bmatrix} \\
&= \begin{bmatrix} 1.0 \\ 0.85 \\ 1.15 \end{bmatrix} + \begin{bmatrix} 0.1 \\ 0.2 \\ 0.1 \end{bmatrix} = \begin{bmatrix} 1.1 \\ 1.05 \\ 1.25 \end{bmatrix}
\end{align}

\textbf{Step 2: Apply ReLU activation}
\begin{equation}
\mathbf{h} = \text{ReLU}(\mathbf{z}) = \max(0, \mathbf{z}) = \begin{bmatrix} 1.1 \\ 1.05 \\ 1.25 \end{bmatrix}
\end{equation}

\textbf{Interpretation:}
\begin{itemize}
    \item Non-square matrix $\mathbf{W} \in \mathbb{R}^{3 \times 4}$ reduces dimension from 4 to 3
    \item Each row of $\mathbf{W}$ computes a weighted sum of inputs
    \item Bias adds learnable offset
    \item ReLU introduces non-linearity
\end{itemize}

\textbf{Training:} Weight matrix $\mathbf{W}$ is learned via backpropagation to minimize loss.
\end{example}

\newpage
\section{Dot Product}

\subsection{Definition and Mathematical Foundation}

\begin{definition}[Dot Product (Inner Product)]
The dot product of two vectors $\mathbf{a}, \mathbf{b} \in \mathbb{R}^n$ is a scalar defined as:
\begin{equation}
\mathbf{a} \cdot \mathbf{b} = \mathbf{a}^T\mathbf{b} = \sum_{i=1}^{n} a_i b_i = a_1b_1 + a_2b_2 + \cdots + a_nb_n
\end{equation}

\textbf{Notation:} Also written as $\langle \mathbf{a}, \mathbf{b} \rangle$ or $(\mathbf{a}, \mathbf{b})$

\textbf{Result:} A single real number (scalar)
\end{definition}

\begin{theorem}[Properties of Dot Product]
For vectors $\mathbf{a}, \mathbf{b}, \mathbf{c} \in \mathbb{R}^n$ and scalar $k \in \mathbb{R}$:

\begin{enumerate}
    \item \textbf{Commutative:} $\mathbf{a} \cdot \mathbf{b} = \mathbf{b} \cdot \mathbf{a}$
    
    \item \textbf{Distributive:} $\mathbf{a} \cdot (\mathbf{b} + \mathbf{c}) = \mathbf{a} \cdot \mathbf{b} + \mathbf{a} \cdot \mathbf{c}$
    
    \item \textbf{Scalar multiplication:} $(k\mathbf{a}) \cdot \mathbf{b} = k(\mathbf{a} \cdot \mathbf{b}) = \mathbf{a} \cdot (k\mathbf{b})$
    
    \item \textbf{Positive definite:} $\mathbf{a} \cdot \mathbf{a} \geq 0$, with equality iff $\mathbf{a} = \mathbf{0}$
    
    \item \textbf{Cauchy-Schwarz inequality:} $|\mathbf{a} \cdot \mathbf{b}| \leq \|\mathbf{a}\| \|\mathbf{b}\|$
\end{enumerate}
\end{theorem}

\subsection{Geometric Interpretation}

\begin{theorem}[Geometric Formula]
The dot product can be expressed using vector magnitudes and the angle between them:
\begin{equation}
\mathbf{a} \cdot \mathbf{b} = \|\mathbf{a}\| \|\mathbf{b}\| \cos \theta
\end{equation}

where $\theta$ is the angle between vectors $\mathbf{a}$ and $\mathbf{b}$.

\textbf{Consequences:}
\begin{itemize}
    \item $\theta = 0°$: $\mathbf{a} \cdot \mathbf{b} = \|\mathbf{a}\| \|\mathbf{b}\|$ (parallel, same direction)
    \item $\theta = 90°$: $\mathbf{a} \cdot \mathbf{b} = 0$ (perpendicular/orthogonal)
    \item $\theta = 180°$: $\mathbf{a} \cdot \mathbf{b} = -\|\mathbf{a}\| \|\mathbf{b}\|$ (parallel, opposite direction)
\end{itemize}
\end{theorem}

\begin{definition}[Vector Magnitude using Dot Product]
The magnitude (or norm) of a vector can be computed using dot product:
\begin{equation}
\|\mathbf{a}\| = \sqrt{\mathbf{a} \cdot \mathbf{a}} = \sqrt{\sum_{i=1}^{n} a_i^2}
\end{equation}

This is also called the Euclidean norm or L2 norm.
\end{definition}

\begin{definition}[Angle Between Vectors]
The angle between two non-zero vectors can be found by rearranging the geometric formula:
\begin{equation}
\cos \theta = \frac{\mathbf{a} \cdot \mathbf{b}}{\|\mathbf{a}\| \|\mathbf{b}\|} \implies \theta = \arccos\left(\frac{\mathbf{a} \cdot \mathbf{b}}{\|\mathbf{a}\| \|\mathbf{b}\|}\right)
\end{equation}
\end{definition}

\subsection{Computational Examples}

\begin{example}[Basic Dot Product Computation]
Compute the dot product of:
\begin{equation}
\mathbf{a} = \begin{bmatrix} 3 \\ 4 \\ 5 \end{bmatrix}, \quad
\mathbf{b} = \begin{bmatrix} 2 \\ -1 \\ 3 \end{bmatrix}
\end{equation}

\textbf{Step 1: Apply formula}
\begin{align}
\mathbf{a} \cdot \mathbf{b} &= \sum_{i=1}^{3} a_i b_i \\
&= (3)(2) + (4)(-1) + (5)(3) \\
&= 6 - 4 + 15 \\
&= 17
\end{align}

\textbf{Step 2: Matrix notation verification}
\begin{equation}
\mathbf{a}^T\mathbf{b} = \begin{bmatrix} 3 & 4 & 5 \end{bmatrix} \begin{bmatrix} 2 \\ -1 \\ 3 \end{bmatrix} = 17 \checkmark
\end{equation}

\textbf{Result:} The dot product is 17 (a positive scalar).
\end{example}

\begin{example}[Computing Angle Between Vectors]
Find the angle between:
\begin{equation}
\mathbf{a} = \begin{bmatrix} 1 \\ 0 \end{bmatrix}, \quad
\mathbf{b} = \begin{bmatrix} 1 \\ 1 \end{bmatrix}
\end{equation}

\textbf{Step 1: Compute dot product}
\begin{equation}
\mathbf{a} \cdot \mathbf{b} = (1)(1) + (0)(1) = 1
\end{equation}

\textbf{Step 2: Compute magnitudes}
\begin{align}
\|\mathbf{a}\| &= \sqrt{1^2 + 0^2} = 1 \\
\|\mathbf{b}\| &= \sqrt{1^2 + 1^2} = \sqrt{2}
\end{align}

\textbf{Step 3: Apply angle formula}
\begin{equation}
\cos \theta = \frac{\mathbf{a} \cdot \mathbf{b}}{\|\mathbf{a}\| \|\mathbf{b}\|} = \frac{1}{1 \cdot \sqrt{2}} = \frac{1}{\sqrt{2}} = \frac{\sqrt{2}}{2}
\end{equation}

\textbf{Step 4: Find angle}
\begin{equation}
\theta = \arccos\left(\frac{\sqrt{2}}{2}\right) = 45° = \frac{\pi}{4} \text{ radians}
\end{equation}

\textbf{Interpretation:} The vectors form a 45° angle, which makes geometric sense: $\mathbf{a}$ points along x-axis, $\mathbf{b}$ points diagonally.
\end{example}

\begin{example}[Testing Orthogonality]
Check if these vectors are orthogonal:
\begin{equation}
\mathbf{a} = \begin{bmatrix} 3 \\ -2 \\ 1 \end{bmatrix}, \quad
\mathbf{b} = \begin{bmatrix} 2 \\ 3 \\ 0 \end{bmatrix}
\end{equation}

\textbf{Computation:}
\begin{equation}
\mathbf{a} \cdot \mathbf{b} = (3)(2) + (-2)(3) + (1)(0) = 6 - 6 + 0 = 0
\end{equation}

\textbf{Conclusion:} Since $\mathbf{a} \cdot \mathbf{b} = 0$, the vectors are orthogonal (perpendicular).

This means $\cos \theta = 0 \implies \theta = 90°$.
\end{example}

\begin{example}[Projection of Vector onto Another]
Project vector $\mathbf{b}$ onto vector $\mathbf{a}$:
\begin{equation}
\mathbf{a} = \begin{bmatrix} 3 \\ 4 \end{bmatrix}, \quad
\mathbf{b} = \begin{bmatrix} 2 \\ 1 \end{bmatrix}
\end{equation}

The projection formula is:
\begin{equation}
\text{proj}_{\mathbf{a}}\mathbf{b} = \frac{\mathbf{a} \cdot \mathbf{b}}{\mathbf{a} \cdot \mathbf{a}} \mathbf{a}
\end{equation}

\textbf{Step 1: Compute $\mathbf{a} \cdot \mathbf{b}$}
\begin{equation}
\mathbf{a} \cdot \mathbf{b} = (3)(2) + (4)(1) = 6 + 4 = 10
\end{equation}

\textbf{Step 2: Compute $\mathbf{a} \cdot \mathbf{a}$}
\begin{equation}
\mathbf{a} \cdot \mathbf{a} = 3^2 + 4^2 = 9 + 16 = 25
\end{equation}

\textbf{Step 3: Calculate projection}
\begin{equation}
\text{proj}_{\mathbf{a}}\mathbf{b} = \frac{10}{25} \begin{bmatrix} 3 \\ 4 \end{bmatrix} = \frac{2}{5} \begin{bmatrix} 3 \\ 4 \end{bmatrix} = \begin{bmatrix} 1.2 \\ 1.6 \end{bmatrix}
\end{equation}

\textbf{Interpretation:} The projection is the component of $\mathbf{b}$ along the direction of $\mathbf{a}$.

\textbf{Verification - projection is parallel to $\mathbf{a}$:}
\begin{equation}
\begin{bmatrix} 1.2 \\ 1.6 \end{bmatrix} = 0.4 \begin{bmatrix} 3 \\ 4 \end{bmatrix} \checkmark
\end{equation}
\end{example}

\subsection{Applications in Machine Learning and Deep Learning}

\begin{tcolorbox}[enhanced,breakable,colback=red!5!white,colframe=red!75!black,title=Critical ML/DL Applications of Dot Product]

\textbf{1. Similarity Measures:}

Cosine similarity between vectors:
\begin{equation}
\text{similarity}(\mathbf{a}, \mathbf{b}) = \frac{\mathbf{a} \cdot \mathbf{b}}{\|\mathbf{a}\| \|\mathbf{b}\|} = \cos \theta
\end{equation}

Range: [-1, 1] where 1 = identical direction, 0 = orthogonal, -1 = opposite

\textbf{Applications:}
\begin{itemize}
    \item Document similarity in NLP
    \item Image similarity in computer vision
    \item Recommendation systems
    \item Clustering algorithms
\end{itemize}

\textbf{2. Neural Network Layers:}

Every neuron computes a weighted sum (dot product):
\begin{equation}
z = \mathbf{w}^T\mathbf{x} + b = \sum_{i=1}^{n} w_i x_i + b
\end{equation}

Then applies activation: $a = \sigma(z)$

\textbf{3. Attention Mechanisms (Transformers):}

Attention scores are computed using dot products:
\begin{equation}
\text{score}(\mathbf{q}, \mathbf{k}) = \frac{\mathbf{q} \cdot \mathbf{k}}{\sqrt{d_k}}
\end{equation}

where $\mathbf{q}$ = query, $\mathbf{k}$ = key, $d_k$ = dimension

\textbf{4. Linear Regression:}

Prediction is a dot product:
\begin{equation}
\hat{y} = \mathbf{w}^T\mathbf{x} = w_1x_1 + w_2x_2 + \cdots + w_dx_d
\end{equation}

\textbf{5. Support Vector Machines (SVMs):}

Decision boundary is based on dot product:
\begin{equation}
f(\mathbf{x}) = \text{sign}(\mathbf{w}^T\mathbf{x} + b)
\end{equation}

Kernel trick: $K(\mathbf{x}_i, \mathbf{x}_j) = \phi(\mathbf{x}_i) \cdot \phi(\mathbf{x}_j)$

\textbf{6. Distance Metrics:}

Euclidean distance uses dot product:
\begin{equation}
\|\mathbf{a} - \mathbf{b}\|^2 = (\mathbf{a} - \mathbf{b}) \cdot (\mathbf{a} - \mathbf{b}) = \mathbf{a} \cdot \mathbf{a} - 2\mathbf{a} \cdot \mathbf{b} + \mathbf{b} \cdot \mathbf{b}
\end{equation}

\textbf{7. Principal Component Analysis:}

Projection onto principal components:
\begin{equation}
z_i = \mathbf{v}_i^T\mathbf{x} \quad \text{(dot product with eigenvector)}
\end{equation}

\textbf{8. Gradient Descent:}

Gradient is computed using dot products in backpropagation:
\begin{equation}
\frac{\partial L}{\partial w_i} = \sum_{j} \frac{\partial L}{\partial z_j} \frac{\partial z_j}{\partial w_i} = \sum_{j} \delta_j x_i
\end{equation}
\end{tcolorbox}

\begin{example}[ML Application: Cosine Similarity for Document Comparison]
Two documents represented as TF-IDF vectors:
\begin{align}
\mathbf{d}_1 &= [0.5, 0.3, 0.0, 0.8, 0.0] \quad \text{(Document 1)} \\
\mathbf{d}_2 &= [0.4, 0.2, 0.1, 0.7, 0.0] \quad \text{(Document 2)}
\end{align}

\textbf{Step 1: Compute dot product}
\begin{align}
\mathbf{d}_1 \cdot \mathbf{d}_2 &= (0.5)(0.4) + (0.3)(0.2) + (0.0)(0.1) + (0.8)(0.7) + (0.0)(0.0) \\
&= 0.20 + 0.06 + 0 + 0.56 + 0 \\
&= 0.82
\end{align}

\textbf{Step 2: Compute magnitudes}
\begin{align}
\|\mathbf{d}_1\| &= \sqrt{0.5^2 + 0.3^2 + 0^2 + 0.8^2 + 0^2} = \sqrt{0.98} \approx 0.990 \\
\|\mathbf{d}_2\| &= \sqrt{0.4^2 + 0.2^2 + 0.1^2 + 0.7^2 + 0^2} = \sqrt{0.70} \approx 0.837
\end{align}

\textbf{Step 3: Compute cosine similarity}
\begin{equation}
\text{similarity} = \frac{\mathbf{d}_1 \cdot \mathbf{d}_2}{\|\mathbf{d}_1\| \|\mathbf{d}_2\|} = \frac{0.82}{0.990 \times 0.837} \approx \frac{0.82}{0.829} \approx 0.989
\end{equation}

\textbf{Interpretation:}
\begin{itemize}
    \item Similarity = 0.989 (very close to 1)
    \item Documents are very similar in content
    \item Would be good candidates for clustering together
\end{itemize}

\textbf{Decision threshold:} If similarity > 0.8, consider documents as duplicates or highly related.
\end{example}

\begin{example}[Deep Learning Application: Neuron Computation]
A single neuron with 4 inputs:

\textbf{Weights:} $\mathbf{w} = [0.5, -0.3, 0.8, 0.2]^T$
\textbf{Bias:} $b = 0.1$
\textbf{Input:} $\mathbf{x} = [1.0, 2.0, 0.5, 1.5]^T$

\textbf{Step 1: Compute weighted sum (dot product)}
\begin{align}
z &= \mathbf{w}^T\mathbf{x} + b \\
&= (0.5)(1.0) + (-0.3)(2.0) + (0.8)(0.5) + (0.2)(1.5) + 0.1 \\
&= 0.5 - 0.6 + 0.4 + 0.3 + 0.1 \\
&= 0.7
\end{align}

\textbf{Step 2: Apply activation function (sigmoid)}
\begin{equation}
a = \sigma(z) = \frac{1}{1 + e^{-z}} = \frac{1}{1 + e^{-0.7}} \approx \frac{1}{1.503} \approx 0.668
\end{equation}

\textbf{Interpretation:}
\begin{itemize}
    \item Neuron receives 4 inputs
    \item Computes weighted sum via dot product
    \item Adds bias term
    \item Applies non-linear activation
    \item Outputs 0.668 (activation level)
\end{itemize}

\textbf{In a neural network:}
\begin{itemize}
    \item This operation happens billions of times during training
    \item Efficient dot product computation is critical for performance
    \item GPUs excel at parallel dot product operations
\end{itemize}
\end{example}

\begin{example}[Transformer Attention: Scaled Dot-Product Attention]
Simplified attention mechanism:

\textbf{Query:} $\mathbf{q} = [1.0, 0.5, 0.8]^T$
\textbf{Keys:} 
\begin{equation}
\mathbf{K} = \begin{bmatrix}
1.2 & 0.3 & 0.7 \\
0.8 & 0.9 & 0.5 \\
0.5 & 0.4 & 1.0
\end{bmatrix}
\end{equation}
(Each row is a key vector)

\textbf{Values:}
\begin{equation}
\mathbf{V} = \begin{bmatrix}
2.0 & 1.5 \\
1.8 & 1.2 \\
1.6 & 2.0
\end{bmatrix}
\end{equation}

\textbf{Step 1: Compute attention scores (dot products)}
\begin{align}
\text{score}_1 &= \mathbf{q} \cdot \mathbf{k}_1 = (1.0)(1.2) + (0.5)(0.3) + (0.8)(0.7) = 1.91 \\
\text{score}_2 &= \mathbf{q} \cdot \mathbf{k}_2 = (1.0)(0.8) + (0.5)(0.9) + (0.8)(0.5) = 1.65 \\
\text{score}_3 &= \mathbf{q} \cdot \mathbf{k}_3 = (1.0)(0.5) + (0.5)(0.4) + (0.8)(1.0) = 1.50
\end{align}

\textbf{Step 2: Scale by $\sqrt{d_k}$}
\begin{equation}
d_k = 3, \quad \sqrt{d_k} = \sqrt{3} \approx 1.732
\end{equation}
\begin{align}
\text{scaled}_1 &= 1.91 / 1.732 \approx 1.103 \\
\text{scaled}_2 &= 1.65 / 1.732 \approx 0.953 \\
\text{scaled}_3 &= 1.50 / 1.732 \approx 0.866
\end{align}

\textbf{Step 3: Apply softmax}
\begin{equation}
e^{1.103} \approx 3.01, \quad e^{0.953} \approx 2.59, \quad e^{0.866} \approx 2.38
\end{equation}
\begin{equation}
\text{sum} = 3.01 + 2.59 + 2.38 = 7.98
\end{equation}
\begin{align}
\alpha_1 &= 3.01 / 7.98 \approx 0.377 \\
\alpha_2 &= 2.59 / 7.98 \approx 0.325 \\
\alpha_3 &= 2.38 / 7.98 \approx 0.298
\end{align}

\textbf{Step 4: Weighted sum of values}
\begin{align}
\text{output} &= \alpha_1\mathbf{v}_1 + \alpha_2\mathbf{v}_2 + \alpha_3\mathbf{v}_3 \\
&= 0.377\begin{bmatrix} 2.0 \\ 1.5 \end{bmatrix} + 0.325\begin{bmatrix} 1.8 \\ 1.2 \end{bmatrix} + 0.298\begin{bmatrix} 1.6 \\ 2.0 \end{bmatrix} \\
&= \begin{bmatrix} 0.754 \\ 0.566 \end{bmatrix} + \begin{bmatrix} 0.585 \\ 0.390 \end{bmatrix} + \begin{bmatrix} 0.477 \\ 0.596 \end{bmatrix} \\
&= \begin{bmatrix} 1.816 \\ 1.552 \end{bmatrix}
\end{align}

\textbf{Interpretation:}
\begin{itemize}
    \item Query attends most to key 1 (highest score)
    \item Attention weights sum to 1
    \item Output is weighted combination of values
    \item Dot products enable efficient similarity computation
\end{itemize}
\end{example}

\newpage
\section{Cross Product}

\subsection{Definition and Mathematical Foundation}

\begin{definition}[Cross Product (Vector Product)]
The cross product of two vectors $\mathbf{a}, \mathbf{b} \in \mathbb{R}^3$ is a vector $\mathbf{c} \in \mathbb{R}^3$ defined as:
\begin{equation}
\mathbf{a} \times \mathbf{b} = \begin{bmatrix}
a_2b_3 - a_3b_2 \\
a_3b_1 - a_1b_3 \\
a_1b_2 - a_2b_1
\end{bmatrix}
\end{equation}

\textbf{Key characteristic:} The result is a vector (unlike dot product which gives a scalar).

\textbf{Restriction:} Only defined for 3-dimensional vectors (or can be extended to 7D, but not commonly used).
\end{definition}

\begin{definition}[Determinant Notation]
The cross product can be expressed as a determinant:
\begin{equation}
\mathbf{a} \times \mathbf{b} = \det\begin{bmatrix}
\mathbf{i} & \mathbf{j} & \mathbf{k} \\
a_1 & a_2 & a_3 \\
b_1 & b_2 & b_3
\end{bmatrix}
\end{equation}

where $\mathbf{i}, \mathbf{j}, \mathbf{k}$ are the standard unit vectors.

\textbf{Expansion:}
\begin{equation}
\mathbf{a} \times \mathbf{b} = (a_2b_3 - a_3b_2)\mathbf{i} - (a_1b_3 - a_3b_1)\mathbf{j} + (a_1b_2 - a_2b_1)\mathbf{k}
\end{equation}
\end{definition}

\begin{theorem}[Properties of Cross Product]
For vectors $\mathbf{a}, \mathbf{b}, \mathbf{c} \in \mathbb{R}^3$ and scalar $k \in \mathbb{R}$:

\begin{enumerate}
    \item \textbf{Anti-commutative:} $\mathbf{a} \times \mathbf{b} = -(\mathbf{b} \times \mathbf{a})$
    
    \item \textbf{Not associative:} $(\mathbf{a} \times \mathbf{b}) \times \mathbf{c} \neq \mathbf{a} \times (\mathbf{b} \times \mathbf{c})$ (in general)
    
    \item \textbf{Distributive:} $\mathbf{a} \times (\mathbf{b} + \mathbf{c}) = \mathbf{a} \times \mathbf{b} + \mathbf{a} \times \mathbf{c}$
    
    \item \textbf{Scalar multiplication:} $(k\mathbf{a}) \times \mathbf{b} = k(\mathbf{a} \times \mathbf{b}) = \mathbf{a} \times (k\mathbf{b})$
    
    \item \textbf{Orthogonality:} $\mathbf{a} \times \mathbf{b}$ is perpendicular to both $\mathbf{a}$ and $\mathbf{b}$:
    \begin{equation}
    (\mathbf{a} \times \mathbf{b}) \cdot \mathbf{a} = 0, \quad (\mathbf{a} \times \mathbf{b}) \cdot \mathbf{b} = 0
    \end{equation}
    
    \item \textbf{Zero product:} $\mathbf{a} \times \mathbf{b} = \mathbf{0} \iff \mathbf{a}$ and $\mathbf{b}$ are parallel
    
    \item \textbf{Self cross product:} $\mathbf{a} \times \mathbf{a} = \mathbf{0}$
\end{enumerate}
\end{theorem}

\subsection{Geometric Interpretation}

\begin{theorem}[Geometric Formula]
The magnitude of the cross product is:
\begin{equation}
\|\mathbf{a} \times \mathbf{b}\| = \|\mathbf{a}\| \|\mathbf{b}\| \sin \theta
\end{equation}

where $\theta$ is the angle between $\mathbf{a}$ and $\mathbf{b}$.

\textbf{Geometric meaning:} $\|\mathbf{a} \times \mathbf{b}\|$ equals the area of the parallelogram formed by $\mathbf{a}$ and $\mathbf{b}$.
\end{theorem}

\begin{theorem}[Right-Hand Rule]
The direction of $\mathbf{a} \times \mathbf{b}$ is determined by the right-hand rule:
\begin{itemize}
    \item Point fingers in direction of $\mathbf{a}$
    \item Curl fingers toward $\mathbf{b}$
    \item Thumb points in direction of $\mathbf{a} \times \mathbf{b}$
\end{itemize}

The resulting vector is perpendicular to the plane containing $\mathbf{a}$ and $\mathbf{b}$.
\end{theorem}

\subsection{Computational Examples}

\begin{example}[Basic Cross Product Computation]
Compute the cross product:
\begin{equation}
\mathbf{a} = \begin{bmatrix} 2 \\ 3 \\ 4 \end{bmatrix}, \quad
\mathbf{b} = \begin{bmatrix} 5 \\ 6 \\ 7 \end{bmatrix}
\end{equation}

\textbf{Method 1: Component formula}
\begin{align}
\mathbf{a} \times \mathbf{b} &= \begin{bmatrix}
a_2b_3 - a_3b_2 \\
a_3b_1 - a_1b_3 \\
a_1b_2 - a_2b_1
\end{bmatrix} \\
&= \begin{bmatrix}
(3)(7) - (4)(6) \\
(4)(5) - (2)(7) \\
(2)(6) - (3)(5)
\end{bmatrix} \\
&= \begin{bmatrix}
21 - 24 \\
20 - 14 \\
12 - 15
\end{bmatrix} = \begin{bmatrix}
-3 \\
6 \\
-3
\end{bmatrix}
\end{align}

\textbf{Method 2: Determinant expansion}
\begin{align}
\mathbf{a} \times \mathbf{b} &= \det\begin{bmatrix}
\mathbf{i} & \mathbf{j} & \mathbf{k} \\
2 & 3 & 4 \\
5 & 6 & 7
\end{bmatrix} \\
&= \mathbf{i}\det\begin{bmatrix} 3 & 4 \\ 6 & 7 \end{bmatrix} - \mathbf{j}\det\begin{bmatrix} 2 & 4 \\ 5 & 7 \end{bmatrix} + \mathbf{k}\det\begin{bmatrix} 2 & 3 \\ 5 & 6 \end{bmatrix} \\
&= \mathbf{i}(21-24) - \mathbf{j}(14-20) + \mathbf{k}(12-15) \\
&= -3\mathbf{i} + 6\mathbf{j} - 3\mathbf{k} = \begin{bmatrix} -3 \\ 6 \\ -3 \end{bmatrix}
\end{align}

\textbf{Verification of orthogonality:}
\begin{align}
(\mathbf{a} \times \mathbf{b}) \cdot \mathbf{a} &= (-3)(2) + (6)(3) + (-3)(4) = -6 + 18 - 12 = 0 \checkmark \\
(\mathbf{a} \times \mathbf{b}) \cdot \mathbf{b} &= (-3)(5) + (6)(6) + (-3)(7) = -15 + 36 - 21 = 0 \checkmark
\end{align}

The cross product is indeed perpendicular to both original vectors!
\end{example}

\begin{example}[Computing Area of Parallelogram]
Find the area of the parallelogram formed by:
\begin{equation}
\mathbf{a} = \begin{bmatrix} 3 \\ 0 \\ 0 \end{bmatrix}, \quad
\mathbf{b} = \begin{bmatrix} 0 \\ 4 \\ 0 \end{bmatrix}
\end{equation}

\textbf{Step 1: Compute cross product}
\begin{align}
\mathbf{a} \times \mathbf{b} &= \begin{bmatrix}
(0)(0) - (0)(4) \\
(0)(0) - (3)(0) \\
(3)(4) - (0)(0)
\end{bmatrix} = \begin{bmatrix}
0 \\
0 \\
12
\end{bmatrix}
\end{align}

\textbf{Step 2: Compute magnitude}
\begin{equation}
\|\mathbf{a} \times \mathbf{b}\| = \sqrt{0^2 + 0^2 + 12^2} = 12
\end{equation}

\textbf{Area:} The area of the parallelogram is 12 square units.

\textbf{Verification using base × height:}
\begin{itemize}
    \item Base = $\|\mathbf{a}\| = 3$
    \item Height = $\|\mathbf{b}\| = 4$
    \item Since vectors are perpendicular ($\theta = 90°$): Area = $3 \times 4 = 12$ ✓
\end{itemize}
\end{example}

\begin{example}[Testing Parallel Vectors]
Check if vectors are parallel:
\begin{equation}
\mathbf{a} = \begin{bmatrix} 2 \\ 4 \\ 6 \end{bmatrix}, \quad
\mathbf{b} = \begin{bmatrix} 1 \\ 2 \\ 3 \end{bmatrix}
\end{equation}

\textbf{Observation:} $\mathbf{a} = 2\mathbf{b}$, so vectors are parallel.

\textbf{Compute cross product:}
\begin{align}
\mathbf{a} \times \mathbf{b} &= \begin{bmatrix}
(4)(3) - (6)(2) \\
(6)(1) - (2)(3) \\
(2)(2) - (4)(1)
\end{bmatrix} \\
&= \begin{bmatrix}
12 - 12 \\
6 - 6 \\
4 - 4
\end{bmatrix} = \begin{bmatrix}
0 \\
0 \\
0
\end{bmatrix}
\end{align}

\textbf{Conclusion:} Cross product is zero vector, confirming vectors are parallel.
\end{example}

\begin{example}[Standard Unit Vector Cross Products]
\textbf{Standard basis vectors:}
\begin{equation}
\mathbf{i} = \begin{bmatrix} 1 \\ 0 \\ 0 \end{bmatrix}, \quad
\mathbf{j} = \begin{bmatrix} 0 \\ 1 \\ 0 \end{bmatrix}, \quad
\mathbf{k} = \begin{bmatrix} 0 \\ 0 \\ 1 \end{bmatrix}
\end{equation}

\textbf{Cyclic pattern (right-hand rule):}
\begin{align}
\mathbf{i} \times \mathbf{j} &= \mathbf{k} \\
\mathbf{j} \times \mathbf{k} &= \mathbf{i} \\
\mathbf{k} \times \mathbf{i} &= \mathbf{j}
\end{align}

\textbf{Anti-cyclic pattern:}
\begin{align}
\mathbf{j} \times \mathbf{i} &= -\mathbf{k} \\
\mathbf{k} \times \mathbf{j} &= -\mathbf{i} \\
\mathbf{i} \times \mathbf{k} &= -\mathbf{j}
\end{align}

\textbf{Self cross products:}
\begin{equation}
\mathbf{i} \times \mathbf{i} = \mathbf{j} \times \mathbf{j} = \mathbf{k} \times \mathbf{k} = \mathbf{0}
\end{equation}

\textbf{Verification - Example:} $\mathbf{i} \times \mathbf{j}$
\begin{equation}
\begin{bmatrix} 1 \\ 0 \\ 0 \end{bmatrix} \times \begin{bmatrix} 0 \\ 1 \\ 0 \end{bmatrix} = \begin{bmatrix}
(0)(0) - (0)(1) \\
(0)(0) - (1)(0) \\
(1)(1) - (0)(0)
\end{bmatrix} = \begin{bmatrix} 0 \\ 0 \\ 1 \end{bmatrix} = \mathbf{k} \checkmark
\end{equation}
\end{example}

\subsection{Applications in Machine Learning and Computer Vision}

\begin{tcolorbox}[enhanced,breakable,colback=purple!5!white,colframe=purple!75!black,title=ML/DL Applications of Cross Product]

\textbf{1. Computer Graphics and 3D Vision:}

\textbf{Surface Normal Computation:}
Given two edges of a triangle:
\begin{equation}
\mathbf{n} = \mathbf{e}_1 \times \mathbf{e}_2
\end{equation}
The cross product gives the normal vector perpendicular to the surface.

\textbf{Applications:}
\begin{itemize}
    \item Lighting calculations (Phong shading)
    \item Back-face culling in rendering
    \item Surface reconstruction
    \item 3D mesh processing
\end{itemize}

\textbf{2. Robotics and Pose Estimation:}

Computing rotational velocities and angular momentum:
\begin{equation}
\boldsymbol{\omega} = \mathbf{r} \times \mathbf{v}
\end{equation}

Used in:
\begin{itemize}
    \item Robot kinematics
    \item Camera pose estimation
    \item SLAM (Simultaneous Localization and Mapping)
\end{itemize}

\textbf{3. Physics-Based Neural Networks:}

Modeling physical constraints:
\begin{itemize}
    \item Torque calculations: $\boldsymbol{\tau} = \mathbf{r} \times \mathbf{F}$
    \item Angular momentum conservation
    \item Rigid body dynamics
\end{itemize}

\textbf{4. Point Cloud Processing:}

Computing surface normals for 3D point clouds:
\begin{equation}
\mathbf{n}_p = (\mathbf{p}_2 - \mathbf{p}_1) \times (\mathbf{p}_3 - \mathbf{p}_1)
\end{equation}

Used in:
\begin{itemize}
    \item 3D object detection
    \item Scene understanding
    \item Autonomous driving
\end{itemize}

\textbf{5. Geometric Deep Learning:}

\begin{itemize}
    \item Graph neural networks with spatial features
    \item Mesh convolutional networks
    \item Point cloud networks (PointNet++)
\end{itemize}

\textbf{6. Camera Calibration:}

Essential matrix and fundamental matrix computations involve cross products for epipolar geometry.

\textbf{7. Volume Computation:}

Scalar triple product (involves cross product):
\begin{equation}
V = |\mathbf{a} \cdot (\mathbf{b} \times \mathbf{c})|
\end{equation}
Volume of parallelepiped formed by three vectors.
\end{tcolorbox}

\begin{example}[Computer Vision: Computing Surface Normal]
A triangular mesh face with vertices:
\begin{align}
\mathbf{p}_1 &= \begin{bmatrix} 1 \\ 0 \\ 0 \end{bmatrix} \\
\mathbf{p}_2 &= \begin{bmatrix} 0 \\ 1 \\ 0 \end{bmatrix} \\
\mathbf{p}_3 &= \begin{bmatrix} 0 \\ 0 \\ 1 \end{bmatrix}
\end{align}

\textbf{Step 1: Compute edge vectors}
\begin{align}
\mathbf{e}_1 &= \mathbf{p}_2 - \mathbf{p}_1 = \begin{bmatrix} -1 \\ 1 \\ 0 \end{bmatrix} \\
\mathbf{e}_2 &= \mathbf{p}_3 - \mathbf{p}_1 = \begin{bmatrix} -1 \\ 0 \\ 1 \end{bmatrix}
\end{align}

\textbf{Step 2: Compute cross product (normal vector)}
\begin{align}
\mathbf{n} &= \mathbf{e}_1 \times \mathbf{e}_2 \\
&= \begin{bmatrix}
(1)(1) - (0)(0) \\
(0)(-1) - (-1)(1) \\
(-1)(0) - (1)(-1)
\end{bmatrix} \\
&= \begin{bmatrix}
1 \\
1 \\
1
\end{bmatrix}
\end{align}

\textbf{Step 3: Normalize to unit normal}
\begin{equation}
\|\mathbf{n}\| = \sqrt{1^2 + 1^2 + 1^2} = \sqrt{3}
\end{equation}
\begin{equation}
\hat{\mathbf{n}} = \frac{\mathbf{n}}{\|\mathbf{n}\|} = \frac{1}{\sqrt{3}}\begin{bmatrix} 1 \\ 1 \\ 1 \end{bmatrix} = \begin{bmatrix} 0.577 \\ 0.577 \\ 0.577 \end{bmatrix}
\end{equation}

\textbf{Application in lighting:}

Lambertian shading intensity:
\begin{equation}
I = k_d (\hat{\mathbf{n}} \cdot \hat{\mathbf{l}})
\end{equation}

where $\hat{\mathbf{l}}$ is the light direction and $k_d$ is the diffuse coefficient.

If light direction is $\hat{\mathbf{l}} = [0, 0, 1]^T$:
\begin{equation}
I = k_d (0.577 \times 0 + 0.577 \times 0 + 0.577 \times 1) = 0.577 k_d
\end{equation}

This determines how bright the surface appears based on its orientation to the light.
\end{example}

\begin{example}[Robotics: Computing Torque]
A robot arm with:
\begin{itemize}
    \item Position vector from joint to force application: $\mathbf{r} = [0.3, 0.2, 0]^T$ meters
    \item Applied force: $\mathbf{F} = [0, 0, -10]^T$ Newtons (downward gravity)
\end{itemize}

\textbf{Compute torque:}
\begin{align}
\boldsymbol{\tau} &= \mathbf{r} \times \mathbf{F} \\
&= \begin{bmatrix} 0.3 \\ 0.2 \\ 0 \end{bmatrix} \times \begin{bmatrix} 0 \\ 0 \\ -10 \end{bmatrix} \\
&= \begin{bmatrix}
(0.2)(-10) - (0)(0) \\
(0)(0) - (0.3)(-10) \\
(0.3)(0) - (0.2)(0)
\end{bmatrix} \\
&= \begin{bmatrix}
-2 \\
3 \\
0
\end{bmatrix} \text{ N·m}
\end{align}

\textbf{Interpretation:}
\begin{itemize}
    \item Torque magnitude: $\|\boldsymbol{\tau}\| = \sqrt{4 + 9} = \sqrt{13} \approx 3.61$ N·m
    \item Direction: perpendicular to both $\mathbf{r}$ and $\mathbf{F}$
    \item This torque causes rotation about the joint
    \item Robot controller must compensate to maintain position
\end{itemize}

\textbf{In ML for robotics:}
\begin{itemize}
    \item Neural networks learn to predict required torques
    \item Physics-informed networks include cross product constraints
    \item Reinforcement learning optimizes control policies
\end{itemize}
\end{example}

\begin{example}[3D Computer Vision: Volume Computation]
Compute volume of tetrahedron with vertices at origin and:
\begin{equation}
\mathbf{a} = \begin{bmatrix} 2 \\ 0 \\ 0 \end{bmatrix}, \quad
\mathbf{b} = \begin{bmatrix} 0 \\ 3 \\ 0 \end{bmatrix}, \quad
\mathbf{c} = \begin{bmatrix} 0 \\ 0 \\ 4 \end{bmatrix}
\end{equation}

\textbf{Formula:} Volume = $\frac{1}{6}|\mathbf{a} \cdot (\mathbf{b} \times \mathbf{c})|$

\textbf{Step 1: Compute cross product}
\begin{align}
\mathbf{b} \times \mathbf{c} &= \begin{bmatrix}
(3)(4) - (0)(0) \\
(0)(0) - (0)(4) \\
(0)(0) - (3)(0)
\end{bmatrix} = \begin{bmatrix}
12 \\
0 \\
0
\end{bmatrix}
\end{align}

\textbf{Step 2: Compute scalar triple product}
\begin{equation}
\mathbf{a} \cdot (\mathbf{b} \times \mathbf{c}) = \begin{bmatrix} 2 \\ 0 \\ 0 \end{bmatrix} \cdot \begin{bmatrix} 12 \\ 0 \\ 0 \end{bmatrix} = 24
\end{equation}

\textbf{Step 3: Compute volume}
\begin{equation}
V = \frac{1}{6}|24| = 4 \text{ cubic units}
\end{equation}

\textbf{Verification:} 
For a tetrahedron with perpendicular edges along axes:
\begin{equation}
V = \frac{1}{6} \times \text{base area} \times \text{height} = \frac{1}{6} \times (2 \times 3) \times 4 = 4 \checkmark
\end{equation}

\textbf{Application:}
\begin{itemize}
    \item Mesh volume computation in 3D reconstruction
    \item Collision detection in physics simulations
    \item Shape analysis in medical imaging
\end{itemize}
\end{example}

\subsection{Cross Product vs Dot Product Summary}

\begin{table}[h]
\centering
\begin{tabular}{|p{3cm}|p{5cm}|p{5cm}|}
\hline
\textbf{Property} & \textbf{Dot Product} & \textbf{Cross Product} \\
\hline
Result & Scalar & Vector \\
\hline
Dimension & Any dimension & Only 3D (primarily) \\
\hline
Formula & $\mathbf{a} \cdot \mathbf{b} = \sum a_i b_i$ & $\mathbf{a} \times \mathbf{b} = $ perpendicular vector \\
\hline
Geometric & $\|\mathbf{a}\|\|\mathbf{b}\|\cos\theta$ & $\|\mathbf{a}\|\|\mathbf{b}\|\sin\theta$ \\
\hline
Commutative? & Yes: $\mathbf{a} \cdot \mathbf{b} = \mathbf{b} \cdot \mathbf{a}$ & No: $\mathbf{a} \times \mathbf{b} = -\mathbf{b} \times \mathbf{a}$ \\
\hline
Measures & Projection, similarity & Area, perpendicularity \\
\hline
Zero when & Perpendicular ($\theta = 90°$) & Parallel ($\theta = 0°$ or $180°$) \\
\hline
ML Uses & Similarity, neurons, attention & Surface normals, 3D geometry \\
\hline
\end{tabular}
\caption{Comparison of Dot Product and Cross Product}
\end{table}

\subsection{Why Cross Product is Less Common in ML}

\begin{tcolorbox}[enhanced,breakable,colback=yellow!10!white,colframe=orange!75!black,title=Limitations of Cross Product in ML]

\textbf{1. Dimension Restriction:}
\begin{itemize}
    \item Only defined for 3D vectors (primarily)
    \item ML often works with high-dimensional data (100s to 1000s of dimensions)
    \item Dot product generalizes to any dimension
\end{itemize}

\textbf{2. Computational Efficiency:}
\begin{itemize}
    \item Dot product: $O(n)$ operations, highly parallelizable
    \item Cross product: More complex, produces vector output
    \item Neural networks prefer scalar outputs from dot products
\end{itemize}

\textbf{3. Application Domains:}
\begin{itemize}
    \item Dot product: Similarity, projections, transformations (universal)
    \item Cross product: Geometric operations, 3D spatial reasoning (specialized)
\end{itemize}

\textbf{4. Where Cross Product Appears:}
\begin{itemize}
    \item Computer graphics and rendering
    \item 3D computer vision (surface normals, pose estimation)
    \item Robotics (torque, angular momentum)
    \item Physics-informed neural networks
    \item Point cloud processing
    \item Geometric deep learning on 3D data
\end{itemize}
\end{tcolorbox}

\section{Conclusion}

This document covered three essential topics for machine learning and deep learning:

\subsection{Key Takeaways}

\begin{enumerate}
    \item \textbf{Linear Transformation of Non-square Matrices:}
    \begin{itemize}
        \item Common in ML due to dimension mismatches
        \item Used in dimensionality reduction (PCA)
        \item Neural network layers often have different input/output sizes
        \item Understanding null space reveals information loss
        \item Pseudo-inverse handles non-invertibility
    \end{itemize}
    
    \item \textbf{Dot Product:}
    \begin{itemize}
        \item Fundamental operation in all ML/DL algorithms
        \item Measures similarity and projection
        \item Core of neural network computations
        \item Used in attention mechanisms (Transformers)
        \item Enables efficient parallel computation
        \item Works in any dimension
    \end{itemize}
    
    \item \textbf{Cross Product:}
    \begin{itemize}
        \item Specialized for 3D geometric operations
        \item Essential in computer vision and graphics
        \item Computes surface normals and areas
        \item Used in robotics for torque calculations
        \item Less common in general ML, but crucial for 3D applications
        \item Produces vectors perpendicular to input vectors
    \end{itemize}
\end{enumerate}

\subsection{Practical Recommendations}

\begin{tcolorbox}[enhanced,breakable,colback=green!5!white,colframe=green!75!black,title=Best Practices]

\textbf{For Implementation:}
\begin{itemize}
    \item Use NumPy/PyTorch for efficient dot product: \texttt{np.dot()}, \texttt{@} operator
    \item Leverage GPU acceleration for matrix operations
    \item Use broadcasting to avoid explicit loops
    \item For cross product: \texttt{np.cross()} in NumPy
    \item Normalize vectors when computing similarity (cosine similarity)
\end{itemize}

\textbf{For ML Applications:}
\begin{itemize}
    \item Dot product: similarity search, neural networks, recommender systems
    \item Cross product: 3D vision, surface reconstruction, robotics
    \item Non-square transformations: dimensionality reduction, feature engineering
    \item Always check dimensions before matrix operations
    \item Use pseudo-inverse for non-square systems
\end{itemize}
\end{tcolorbox}

\subsection{Further Study}

To deepen understanding:
\begin{itemize}
    \item Study tensor operations for modern deep learning
    \item Learn attention mechanisms in detail (Transformers)
    \item Explore 3D deep learning (PointNet, Mesh R-CNN)
    \item Understand matrix calculus for backpropagation
    \item Practice implementing algorithms from scratch
    \item Study geometric deep learning frameworks
\end{itemize}

\subsection{Connection to Previous Topics}

These concepts build upon:
\begin{itemize}
    \item \textbf{Matrix inverse:} Pseudo-inverse for non-square matrices
    \item \textbf{Rank:} Determines dimension of transformed space
    \item \textbf{Null space:} Information lost in non-square transformations
    \item \textbf{Column space:} Range of transformation
    \item \textbf{Change of basis:} PCA as change to principal component basis
\end{itemize}

\newpage

\section{Assessment Questions}

\subsection{Level 1: Basic Understanding}

\begin{enumerate}
\item \textbf{What is a non-square matrix?} Provide two examples from machine learning.

\item \textbf{Give the definition} of dot product and cross product. What is the key difference in their outputs?

\item \textbf{What is the difference between} a left inverse and a right inverse? When does each exist?

\item \textbf{Calculate}: Compute the dot product of $\mathbf{a} = \begin{bmatrix} 2 \\ 3 \\ 1 \end{bmatrix}$ and $\mathbf{b} = \begin{bmatrix} 1 \\ -1 \\ 2 \end{bmatrix}$.

\item \textbf{If} $\mathbf{A} = \begin{bmatrix} 1 & 2 \\ 0 & 1 \\ 1 & 0 \end{bmatrix}$ and $\mathbf{x} = \begin{bmatrix} 2 \\ 3 \end{bmatrix}$, find $\mathbf{Ax}$.

\item \textbf{True or False}: The dot product of two vectors is always a scalar.

\item \textbf{Identify whether each operation} produces a scalar or vector:
\begin{enumerate}[label=(\alph*)]
\item Dot product: $\mathbf{a} \cdot \mathbf{b}$
\item Cross product: $\mathbf{a} \times \mathbf{b}$
\item Matrix-vector product: $\mathbf{A}\mathbf{x}$
\end{enumerate}

\item \textbf{Compute}: $\mathbf{a} \times \mathbf{b}$ where $\mathbf{a} = \begin{bmatrix} 1 \\ 0 \\ 0 \end{bmatrix}$ and $\mathbf{b} = \begin{bmatrix} 0 \\ 1 \\ 0 \end{bmatrix}$.

\item \textbf{What are the two fundamental properties} that define a linear transformation?

\item \textbf{Why is the cross product} only defined for 3D vectors (primarily)?
\end{enumerate}

\subsection{Level 2: Intermediate}

\begin{enumerate}[resume]
\item \textbf{What is the role of dot product} in neural network layers? Provide the formula.

\item \textbf{Compute} the left inverse of $\mathbf{A} = \begin{bmatrix} 1 & 0 \\ 0 & 1 \\ 1 & 1 \end{bmatrix}$.

\item \textbf{State the geometric formula} for dot product and explain its meaning.

\item \textbf{Given} $\mathbf{a} = \begin{bmatrix} 3 \\ 4 \end{bmatrix}$ and $\mathbf{b} = \begin{bmatrix} 2 \\ 1 \end{bmatrix}$, compute the projection of $\mathbf{b}$ onto $\mathbf{a}$.

\item \textbf{What is the null space} of a non-square matrix? Give geometric interpretation.

\item \textbf{Given transformation} matrix $\mathbf{A} = \begin{bmatrix} 1 & 2 \\ 0 & 1 \\ 1 & 0 \end{bmatrix}$, find the dimension of its column space and null space.

\item \textbf{Calculate}: Find the angle between $\mathbf{a} = \begin{bmatrix} 1 \\ 0 \end{bmatrix}$ and $\mathbf{b} = \begin{bmatrix} 1 \\ 1 \end{bmatrix}$ using dot product.

\item \textbf{If} $\mathbf{a} = \begin{bmatrix} 2 \\ 3 \\ 4 \end{bmatrix}$ and $\mathbf{b} = \begin{bmatrix} 5 \\ 6 \\ 7 \end{bmatrix}$, compute $\mathbf{a} \times \mathbf{b}$.

\item \textbf{Solve}: For $\mathbf{A} \in \mathbb{R}^{3 \times 2}$ with full column rank, find the left inverse formula.

\item \textbf{Calculate} the area of parallelogram formed by $\mathbf{a} = \begin{bmatrix} 3 \\ 0 \\ 0 \end{bmatrix}$ and $\mathbf{b} = \begin{bmatrix} 0 \\ 4 \\ 0 \end{bmatrix}$.
\end{enumerate}

\subsection{Level 3: Advanced}

\begin{enumerate}[resume]
\item \textbf{Prove} that the dot product is commutative: $\mathbf{a} \cdot \mathbf{b} = \mathbf{b} \cdot \mathbf{a}$.

\item \textbf{Why does a non-square matrix} $\mathbf{A} \in \mathbb{R}^{m \times n}$ with $m > n$ have a left inverse but not a right inverse?

\item \textbf{Explain the role of dot product} in attention mechanisms (Transformers). Provide the scaled dot-product attention formula.

\item \textbf{Compare the computational complexity} of computing:
\begin{enumerate}[label=(\alph*)]
\item Dot product of two $n$-dimensional vectors
\item Cross product of two 3D vectors
\item Matrix-vector product $\mathbf{A}\mathbf{x}$ where $\mathbf{A} \in \mathbb{R}^{m \times n}$
\end{enumerate}

\item \textbf{Design a PCA transformation} using non-square matrix to reduce 1000-dimensional data to 50 dimensions. Specify the transformation matrix dimensions.

\item \textbf{Explain how cross product} is used in computing surface normals for 3D computer vision. Provide the formula.

\item \textbf{What happens to the rank} when you apply a non-square transformation $\mathbf{A} \in \mathbb{R}^{m \times n}$ where $m < n$? Prove your answer.

\item \textbf{Prove} that the cross product is anti-commutative: $\mathbf{a} \times \mathbf{b} = -(\mathbf{b} \times \mathbf{a})$.

\item \textbf{Explain the relationship between} null space dimension and information loss in dimensionality reduction.

\item \textbf{How do you determine} if a non-square matrix has full column rank or full row rank? What are the implications for ML applications?
\end{enumerate}

\subsection{Multiple Choice Questions}

\begin{enumerate}[resume]
\item The dot product of two vectors $\mathbf{a}, \mathbf{b} \in \mathbb{R}^n$ is:
\begin{enumerate}[label=(\alph*)]
\item A vector in $\mathbb{R}^n$
\item A scalar
\item A matrix in $\mathbb{R}^{n \times n}$
\item A tensor
\end{enumerate}

\item For a non-square matrix $\mathbf{A} \in \mathbb{R}^{m \times n}$ with $m > n$ and full column rank:
\begin{enumerate}[label=(\alph*)]
\item Only right inverse exists
\item Only left inverse exists
\item Both inverses exist
\item No inverse exists
\end{enumerate}

\item The cross product $\mathbf{a} \times \mathbf{b}$ is zero if and only if:
\begin{enumerate}[label=(\alph*)]
\item $\mathbf{a}$ and $\mathbf{b}$ are perpendicular
\item $\mathbf{a}$ and $\mathbf{b}$ are parallel
\item $\mathbf{a}$ or $\mathbf{b}$ is zero
\item Both (b) and (c)
\end{enumerate}

\item In a neural network layer with input dimension 784 and output dimension 128, the weight matrix $\mathbf{W}$ has dimensions:
\begin{enumerate}[label=(\alph*)]
\item $784 \times 128$
\item $128 \times 784$
\item $784 \times 784$
\item $128 \times 128$
\end{enumerate}

\item The computational complexity of dot product for two $n$-dimensional vectors is:
\begin{enumerate}[label=(\alph*)]
\item $O(n^2)$
\item $O(n)$
\item $O(1)$
\item $O(\log n)$
\end{enumerate}
\end{enumerate}

\subsection{Programming Exercises}

\begin{enumerate}[label=\textbf{P\arabic*.}]
\item Implement dot product computation in Python without using NumPy's built-in functions.

\item Write a function to compute the left inverse of a non-square matrix $\mathbf{A} \in \mathbb{R}^{m \times n}$ where $m > n$.

\item Create a function to compute cross product of two 3D vectors and verify orthogonality.

\item Implement cosine similarity between two document vectors using dot product.

\item Build a simple PCA dimensionality reduction function that uses non-square transformation matrix.
\end{enumerate}

\newpage

\section{Solutions to Assessment Questions}

\subsection{Level 1 Solutions}

\textbf{A1.} A non-square matrix is a matrix $\mathbf{A} \in \mathbb{R}^{m \times n}$ where $m \neq n$. 

\textit{ML Examples:}
\begin{enumerate}
\item \textit{Neural Network Layer}: Weight matrix $\mathbf{W} \in \mathbb{R}^{128 \times 784}$ maps 784-dimensional input to 128-dimensional output.
\item \textit{PCA Transformation}: Projection matrix $\mathbf{W} \in \mathbb{R}^{50 \times 1000}$ reduces 1000-dimensional features to 50 dimensions.
\end{enumerate}

\textit{Explanation}: Non-square matrices are common in ML because input and output dimensions often differ.

\textbf{A2.} \textbf{Dot Product}: $\mathbf{a} \cdot \mathbf{b} = \sum_{i=1}^{n} a_i b_i = \mathbf{a}^T\mathbf{b}$ (produces a scalar)

\textbf{Cross Product}: $\mathbf{a} \times \mathbf{b}$ produces a vector perpendicular to both $\mathbf{a}$ and $\mathbf{b}$ (only for 3D vectors)

\textit{Key difference}: Dot product gives a scalar; cross product gives a vector.

\textbf{A3.} \textbf{Left Inverse}: For $\mathbf{A} \in \mathbb{R}^{m \times n}$ with $m > n$ and full column rank, $\mathbf{A}_L^{-1} = (\mathbf{A}^T\mathbf{A})^{-1}\mathbf{A}^T$ where $\mathbf{A}_L^{-1}\mathbf{A} = \mathbf{I}_n$.

\textbf{Right Inverse}: For $\mathbf{A} \in \mathbb{R}^{m \times n}$ with $m < n$ and full row rank, $\mathbf{A}_R^{-1} = \mathbf{A}^T(\mathbf{A}\mathbf{A}^T)^{-1}$ where $\mathbf{A}\mathbf{A}_R^{-1} = \mathbf{I}_m$.

\textit{Explanation}: Left inverse exists when there are more rows than columns; right inverse exists when there are more columns than rows.

\textbf{A4.}
\begin{align}
\mathbf{a} \cdot \mathbf{b} &= (2)(1) + (3)(-1) + (1)(2) \\
&= 2 - 3 + 2 \\
&= 1
\end{align}

\textbf{A5.}
\begin{align}
\mathbf{Ax} &= \begin{bmatrix} 1 & 2 \\ 0 & 1 \\ 1 & 0 \end{bmatrix} \begin{bmatrix} 2 \\ 3 \end{bmatrix} \\
&= \begin{bmatrix} 1(2) + 2(3) \\ 0(2) + 1(3) \\ 1(2) + 0(3) \end{bmatrix} \\
&= \begin{bmatrix} 8 \\ 3 \\ 2 \end{bmatrix}
\end{align}

\textbf{A6.} \textbf{True}. The dot product always produces a single scalar value.

\textit{Explanation}: $\mathbf{a} \cdot \mathbf{b} = \sum_{i=1}^{n} a_i b_i$ is a sum of products, resulting in one number.

\textbf{A7.}
\begin{enumerate}[label=(\alph*)]
\item Dot product: \textbf{Scalar}
\item Cross product: \textbf{Vector}
\item Matrix-vector product: \textbf{Vector}
\end{enumerate}

\textbf{A8.}
\begin{align}
\mathbf{a} \times \mathbf{b} &= \begin{bmatrix} 1 \\ 0 \\ 0 \end{bmatrix} \times \begin{bmatrix} 0 \\ 1 \\ 0 \end{bmatrix} \\
&= \begin{bmatrix}
(0)(0) - (0)(1) \\
(0)(0) - (1)(0) \\
(1)(1) - (0)(0)
\end{bmatrix} \\
&= \begin{bmatrix} 0 \\ 0 \\ 1 \end{bmatrix} = \mathbf{k}
\end{align}

\textit{Explanation}: This follows the right-hand rule: $\mathbf{i} \times \mathbf{j} = \mathbf{k}$.

\textbf{A9.} The two fundamental properties are:
\begin{enumerate}
\item \textbf{Additivity}: $T(\mathbf{u} + \mathbf{v}) = T(\mathbf{u}) + T(\mathbf{v})$
\item \textbf{Homogeneity}: $T(c\mathbf{u}) = cT(\mathbf{u})$ for any scalar $c$
\end{enumerate}

\textbf{A10.} The cross product is primarily defined for 3D vectors because:
\begin{itemize}
\item The formula requires exactly 3 components to compute the perpendicular vector
\item The right-hand rule applies specifically to 3D space
\item While it can be extended to 7D, it's not commonly used in higher dimensions
\item Most geometric applications (surface normals, torque) are in 3D space
\end{itemize}

\subsection{Level 2 Solutions}

\textbf{A11.} In neural network layers, dot product computes the weighted sum:
\[
z = \mathbf{w}^T\mathbf{x} + b = \sum_{i=1}^{n} w_i x_i + b
\]

\textit{Explanation}: Each neuron computes a dot product between input vector and weight vector, then adds bias and applies activation function.

\textbf{A12.} For $\mathbf{A} = \begin{bmatrix} 1 & 0 \\ 0 & 1 \\ 1 & 1 \end{bmatrix}$:

\textit{Step 1}: Compute $\mathbf{A}^T\mathbf{A}$
\[
\mathbf{A}^T\mathbf{A} = \begin{bmatrix} 1 & 0 & 1 \\ 0 & 1 & 1 \end{bmatrix} \begin{bmatrix} 1 & 0 \\ 0 & 1 \\ 1 & 1 \end{bmatrix} = \begin{bmatrix} 2 & 1 \\ 1 & 2 \end{bmatrix}
\]

\textit{Step 2}: Invert $\mathbf{A}^T\mathbf{A}$
\[
(\mathbf{A}^T\mathbf{A})^{-1} = \frac{1}{3}\begin{bmatrix} 2 & -1 \\ -1 & 2 \end{bmatrix}
\]

\textit{Step 3}: Compute left inverse
\[
\mathbf{A}_L^{-1} = \frac{1}{3}\begin{bmatrix} 2 & -1 \\ -1 & 2 \end{bmatrix} \begin{bmatrix} 1 & 0 & 1 \\ 0 & 1 & 1 \end{bmatrix} = \frac{1}{3}\begin{bmatrix} 2 & -1 & 1 \\ -1 & 2 & 1 \end{bmatrix}
\]

\textbf{A13.} \textbf{Geometric formula}: $\mathbf{a} \cdot \mathbf{b} = \|\mathbf{a}\| \|\mathbf{b}\| \cos \theta$

\textit{Meaning}: The dot product equals the product of vector magnitudes times the cosine of the angle between them. This relates algebraic computation to geometric relationships.

\textbf{A14.} Projection formula: $\text{proj}_{\mathbf{a}}\mathbf{b} = \frac{\mathbf{a} \cdot \mathbf{b}}{\mathbf{a} \cdot \mathbf{a}} \mathbf{a}$

\textit{Step 1}: $\mathbf{a} \cdot \mathbf{b} = (3)(2) + (4)(1) = 10$

\textit{Step 2}: $\mathbf{a} \cdot \mathbf{a} = 3^2 + 4^2 = 25$

\textit{Step 3}: $\text{proj}_{\mathbf{a}}\mathbf{b} = \frac{10}{25} \begin{bmatrix} 3 \\ 4 \end{bmatrix} = \begin{bmatrix} 1.2 \\ 1.6 \end{bmatrix}$

\textbf{A15.} The \textbf{null space} of matrix $\mathbf{A}$ is:
\[
\mathcal{N}(\mathbf{A}) = \{\mathbf{x} : \mathbf{A}\mathbf{x} = \mathbf{0}\}
\]

\textit{Geometric interpretation}: All vectors that get mapped to zero by the transformation. For non-square matrices with $m < n$, the null space is non-trivial, meaning information is lost.

\textit{Example}: In dimensionality reduction, vectors in the null space direction are collapsed to zero.

\textbf{A16.} For $\mathbf{A} = \begin{bmatrix} 1 & 2 \\ 0 & 1 \\ 1 & 0 \end{bmatrix} \in \mathbb{R}^{3 \times 2}$:

\begin{itemize}
\item $\text{rank}(\mathbf{A}) = 2$ (columns are linearly independent)
\item $\dim(\mathcal{C}(\mathbf{A})) = 2$ (column space dimension equals rank)
\item $\dim(\mathcal{N}(\mathbf{A})) = 2 - 2 = 0$ (null space is trivial since $m > n$)
\end{itemize}

\textbf{A17.} Using dot product formula:
\begin{align}
\mathbf{a} \cdot \mathbf{b} &= (1)(1) + (0)(1) = 1 \\
\|\mathbf{a}\| &= 1, \quad \|\mathbf{b}\| = \sqrt{2} \\
\cos \theta &= \frac{1}{1 \cdot \sqrt{2}} = \frac{1}{\sqrt{2}} \\
\theta &= \arccos\left(\frac{1}{\sqrt{2}}\right) = 45° = \frac{\pi}{4}
\end{align}

\textbf{A18.}
\begin{align}
\mathbf{a} \times \mathbf{b} &= \begin{bmatrix}
(3)(7) - (4)(6) \\
(4)(5) - (2)(7) \\
(2)(6) - (3)(5)
\end{bmatrix} \\
&= \begin{bmatrix}
21 - 24 \\
20 - 14 \\
12 - 15
\end{bmatrix} = \begin{bmatrix} -3 \\ 6 \\ -3 \end{bmatrix}
\end{align}

\textbf{A19.} For $\mathbf{A} \in \mathbb{R}^{m \times n}$ with $m > n$ and full column rank:

\textbf{Left inverse formula}:
\[
\mathbf{A}_L^{-1} = (\mathbf{A}^T\mathbf{A})^{-1}\mathbf{A}^T
\]

\textit{Verification}: $\mathbf{A}_L^{-1}\mathbf{A} = (\mathbf{A}^T\mathbf{A})^{-1}\mathbf{A}^T\mathbf{A} = \mathbf{I}_n$ $\checkmark$

\textbf{A20.} Area of parallelogram = $\|\mathbf{a} \times \mathbf{b}\|$

\begin{align}
\mathbf{a} \times \mathbf{b} &= \begin{bmatrix}
(0)(0) - (0)(4) \\
(0)(0) - (3)(0) \\
(3)(4) - (0)(0)
\end{bmatrix} = \begin{bmatrix} 0 \\ 0 \\ 12 \end{bmatrix} \\
\|\mathbf{a} \times \mathbf{b}\| &= 12
\end{align}

\textit{Verification}: Since vectors are perpendicular, area = $3 \times 4 = 12$ $\checkmark$

\subsection{Level 3 Solutions}

\textbf{A21.} \textbf{Proof of commutativity}:

Let $\mathbf{a} = [a_1, a_2, \ldots, a_n]^T$ and $\mathbf{b} = [b_1, b_2, \ldots, b_n]^T$.

\begin{align}
\mathbf{a} \cdot \mathbf{b} &= \sum_{i=1}^{n} a_i b_i \\
&= \sum_{i=1}^{n} b_i a_i \quad \text{(multiplication is commutative)} \\
&= \mathbf{b} \cdot \mathbf{a}
\end{align}

Therefore, $\mathbf{a} \cdot \mathbf{b} = \mathbf{b} \cdot \mathbf{a}$ $\checkmark$

\textbf{A22.} For $\mathbf{A} \in \mathbb{R}^{m \times n}$ with $m > n$:

\begin{itemize}
\item \textbf{Left inverse exists} if $\text{rank}(\mathbf{A}) = n$ (full column rank)
\item \textbf{Right inverse cannot exist} because:
  \begin{itemize}
  \item Right inverse requires $\mathbf{A}\mathbf{A}_R^{-1} = \mathbf{I}_m$
  \item But $\mathbf{A}$ has only $n$ columns, so $\mathbf{A}\mathbf{A}_R^{-1}$ has at most rank $n < m$
  \item Cannot equal $\mathbf{I}_m$ which has rank $m$
  \end{itemize}
\end{itemize}

\textit{Geometric interpretation}: With more rows than columns, the transformation maps to a lower-dimensional space, so we can recover the input (left inverse) but cannot create a full identity on the output space.

\textbf{A23.} In attention mechanisms, dot product computes similarity scores:

\textbf{Scaled Dot-Product Attention}:
\[
\text{Attention}(\mathbf{Q}, \mathbf{K}, \mathbf{V}) = \text{softmax}\left(\frac{\mathbf{Q}\mathbf{K}^T}{\sqrt{d_k}}\right)\mathbf{V}
\]

\textit{Key role}:
\begin{itemize}
\item $\mathbf{Q}\mathbf{K}^T$ computes dot products between queries and keys
\item Higher dot product = higher similarity = more attention
\item Scaling by $\sqrt{d_k}$ prevents extreme values
\item Softmax converts scores to probability distribution
\end{itemize}

\textbf{A24.} \textbf{Computational complexity}:

\textbf{(a) Dot product}: $O(n)$ operations (n multiplications + n-1 additions)

\textbf{(b) Cross product}: $O(1)$ operations (fixed 9 operations for 3D vectors)

\textbf{(c) Matrix-vector product}: $O(mn)$ operations (m dot products, each O(n))

\textit{Summary}: Dot product is linear, cross product is constant (for 3D), matrix-vector product is quadratic in dimensions.

\textbf{A25.} \textbf{PCA Transformation Design}:

\textit{Transformation}: $\mathbf{Z} = \mathbf{X}\mathbf{W}$ where:
\begin{itemize}
\item $\mathbf{X} \in \mathbb{R}^{N \times 1000}$ (N samples, 1000 features)
\item $\mathbf{W} \in \mathbb{R}^{1000 \times 50}$ (non-square transformation matrix)
\item $\mathbf{Z} \in \mathbb{R}^{N \times 50}$ (reduced dimension)
\end{itemize}

\textit{Steps}:
\begin{enumerate}
\item Center data: $\mathbf{X}_{centered} = \mathbf{X} - \boldsymbol{\mu}$
\item Compute covariance: $\mathbf{C} = \frac{1}{N-1}\mathbf{X}_{centered}^T\mathbf{X}_{centered}$
\item Eigendecomposition: $\mathbf{C} = \mathbf{V}\boldsymbol{\Lambda}\mathbf{V}^T$
\item Select top 50 eigenvectors: $\mathbf{W} = \mathbf{V}_{:,:50}$ (first 50 columns)
\item Project: $\mathbf{Z} = \mathbf{X}_{centered}\mathbf{W}$
\end{enumerate}

\textbf{A26.} \textbf{Surface Normal Computation}:

For triangular mesh with vertices $\mathbf{p}_1, \mathbf{p}_2, \mathbf{p}_3$:

\begin{align}
\mathbf{e}_1 &= \mathbf{p}_2 - \mathbf{p}_1 \\
\mathbf{e}_2 &= \mathbf{p}_3 - \mathbf{p}_1 \\
\mathbf{n} &= \mathbf{e}_1 \times \mathbf{e}_2 \quad \text{(normal vector)} \\
\hat{\mathbf{n}} &= \frac{\mathbf{n}}{\|\mathbf{n}\|} \quad \text{(unit normal)}
\end{align}

\textit{Application}: Used in lighting calculations, back-face culling, and 3D reconstruction.

\textbf{A27.} \textbf{Rank under non-square transformation}:

For $\mathbf{A} \in \mathbb{R}^{m \times n}$ with $m < n$:

\textit{Theorem}: $\text{rank}(\mathbf{A}) \leq \min(m, n) = m$

\textit{Proof}: 
\begin{itemize}
\item Rank equals dimension of column space
\item Column space is subspace of $\mathbb{R}^m$
\item Maximum dimension of subspace in $\mathbb{R}^m$ is $m$
\item Therefore, $\text{rank}(\mathbf{A}) \leq m$
\end{itemize}

\textit{Implication}: When $m < n$, we cannot have full column rank, so information is lost (null space is non-trivial).

\textbf{A28.} \textbf{Proof of anti-commutativity}:

\begin{align}
\mathbf{a} \times \mathbf{b} &= \begin{bmatrix}
a_2b_3 - a_3b_2 \\
a_3b_1 - a_1b_3 \\
a_1b_2 - a_2b_1
\end{bmatrix} \\
\mathbf{b} \times \mathbf{a} &= \begin{bmatrix}
b_2a_3 - b_3a_2 \\
b_3a_1 - b_1a_3 \\
b_1a_2 - b_2a_1
\end{bmatrix} \\
&= \begin{bmatrix}
-(a_2b_3 - a_3b_2) \\
-(a_3b_1 - a_1b_3) \\
-(a_1b_2 - a_2b_1)
\end{bmatrix} \\
&= -(\mathbf{a} \times \mathbf{b})
\end{align}

Therefore, $\mathbf{a} \times \mathbf{b} = -(\mathbf{b} \times \mathbf{a})$ $\checkmark$

\textbf{A29.} \textbf{Null space and information loss}:

\textit{Relationship}: $\dim(\mathcal{N}(\mathbf{A})) = n - \text{rank}(\mathbf{A})$

\textit{In dimensionality reduction}:
\begin{itemize}
\item Larger null space $\implies$ more information lost
\item Vectors in null space are mapped to zero
\item Only information in column space is preserved
\item Rank-nullity theorem: $\dim(\mathcal{C}(\mathbf{A})) + \dim(\mathcal{N}(\mathbf{A})) = n$
\end{itemize}

\textit{Example}: PCA with $\mathbf{W} \in \mathbb{R}^{1000 \times 50}$:
\begin{itemize}
\item Null space dimension = $1000 - 50 = 950$
\item 950 dimensions of information are lost
\item Only 50 principal directions are preserved
\end{itemize}

\textbf{A30.} \textbf{Determining rank properties}:

For $\mathbf{A} \in \mathbb{R}^{m \times n}$:

\textbf{Full column rank} ($\text{rank}(\mathbf{A}) = n$):
\begin{itemize}
\item All columns are linearly independent
\item Null space is trivial: $\dim(\mathcal{N}(\mathbf{A})) = 0$
\item Left inverse exists if $m > n$
\item No information loss in transformation
\end{itemize}

\textbf{Full row rank} ($\text{rank}(\mathbf{A}) = m$):
\begin{itemize}
\item All rows are linearly independent
\item Column space spans entire $\mathbb{R}^m$
\item Right inverse exists if $m < n$
\item Transformation is surjective
\end{itemize}

\textit{ML implications}:
\begin{itemize}
\item Full column rank: Good for feature extraction (no redundancy)
\item Full row rank: Good for covering output space completely
\item Rank deficiency: Indicates redundant features or insufficient model capacity
\end{itemize}

\subsection{Multiple Choice Solutions}

\textbf{A31.} \textbf{(b) A scalar}

\textit{Explanation}: The dot product $\mathbf{a} \cdot \mathbf{b} = \sum_{i=1}^{n} a_i b_i$ is a sum of products, resulting in a single real number (scalar).

\textbf{A32.} \textbf{(b) Only left inverse exists}

\textit{Explanation}: For $m > n$ with full column rank, $\mathbf{A}^T\mathbf{A}$ is invertible, allowing left inverse. Right inverse cannot exist because $\mathbf{A}$ cannot span full $\mathbb{R}^m$ space.

\textbf{A33.} \textbf{(d) Both (b) and (c)}

\textit{Explanation}: Cross product is zero when vectors are parallel ($\mathbf{a} = k\mathbf{b}$) OR when either vector is zero. Both conditions make $\sin \theta = 0$.

\textbf{A34.} \textbf{(b) $128 \times 784$}

\textit{Explanation}: For transformation $\mathbf{h} = \mathbf{W}\mathbf{x} + \mathbf{b}$ where $\mathbf{x} \in \mathbb{R}^{784}$ and $\mathbf{h} \in \mathbb{R}^{128}$, we need $\mathbf{W} \in \mathbb{R}^{128 \times 784}$ to satisfy dimension compatibility: $(128 \times 784)(784 \times 1) = (128 \times 1)$.

\textbf{A35.} \textbf{(b) $O(n)$}

\textit{Explanation}: Dot product requires $n$ multiplications and $n-1$ additions, giving $O(n)$ time complexity. This is linear in the vector dimension.

\newpage

\section{Practice Problems}

\subsection{Computational Problems}

\begin{enumerate}
\item \textbf{Calculate}: Compute dot product of $\mathbf{a} = [1, 2, 3, 4]^T$ and $\mathbf{b} = [2, -1, 0, 1]^T$.

\item \textbf{If} $\mathbf{A} = \begin{bmatrix} 2 & 1 \\ 0 & 3 \\ 1 & 2 \end{bmatrix}$, find $\mathbf{A}\mathbf{x}$ where $\mathbf{x} = \begin{bmatrix} 1 \\ 2 \end{bmatrix}$.

\item \textbf{Compute} cross product: $\mathbf{a} = \begin{bmatrix} 1 \\ 2 \\ 3 \end{bmatrix}$ and $\mathbf{b} = \begin{bmatrix} 4 \\ 5 \\ 6 \end{bmatrix}$.

\item \textbf{Calculate} the angle between $\mathbf{a} = \begin{bmatrix} 1 \\ 1 \\ 0 \end{bmatrix}$ and $\mathbf{b} = \begin{bmatrix} 0 \\ 1 \\ 1 \end{bmatrix}$.

\item \textbf{Find} the projection of $\mathbf{b} = \begin{bmatrix} 2 \\ 1 \end{bmatrix}$ onto $\mathbf{a} = \begin{bmatrix} 3 \\ 4 \end{bmatrix}$.
\end{enumerate}

\subsection{Conceptual Problems}

\begin{enumerate}[resume]
\item \textbf{Why can't} a non-square matrix $\mathbf{A} \in \mathbb{R}^{m \times n}$ with $m \neq n$ have a true inverse?

\item \textbf{How does} the dot product relate to cosine similarity in machine learning?

\item \textbf{Explain} why the cross product is perpendicular to both input vectors.

\item \textbf{Why does} PCA use a non-square transformation matrix?

\item \textbf{How do} dot products enable efficient computation in neural networks?
\end{enumerate}

\subsection{Application Problems}

\begin{enumerate}[resume]
\item \textbf{Design} a neural network layer that transforms 256-dimensional input to 128-dimensional output. Specify the weight matrix dimensions.

\item \textbf{Implement} cosine similarity computation between two document vectors.

\item \textbf{Create} a function to compute surface normal for a triangular face in 3D space.

\item \textbf{Calculate} the left inverse for a given non-square matrix with full column rank.

\item \textbf{Derive} the formula for projection of a vector onto another using dot product.
\end{enumerate}

\vspace{1cm}
\begin{center}
\textbf{End of Notes}

\vspace{0.5cm}
These comprehensive notes provide the mathematical foundation and practical applications needed for understanding modern machine learning and deep learning algorithms.
\end{center}

\end{document}